\documentclass[11pt]{article}
\usepackage[utf8]{inputenc}
\usepackage[colorlinks=true, linkcolor=blue, urlcolor=blue]{hyperref}
\usepackage{geometry}
\geometry{a4paper, margin=1in}
\usepackage{enumitem}

\title{Project Bellwether: Prioritized Intelligence Dossier}
\author{Interview Preparation: Lead Machine Learning Engineer}
\date{March 18, 2026}

\begin{document}
\maketitle

\noindent\textit{This document is rigorously vetted via multi-source verification (excluding single-source LinkedIn where blocked) and strictly ordered from highest immediate relevance (today's recruiter call) to broader context (late-stage interviews).}

\section{Priority 1: The Talent Pipeline (Your Immediate Audience)}
Your call today is with Gabe Schook (Talent Sourcer). He operates within the broader X Talent ecosystem. Understanding his superiors gives you leverage on how the hiring pipeline is structured.
\begin{itemize}[leftmargin=*]
    \item \textbf{Supreetha Ravikumar:} Talent Lead at X (via Talentful). Likely manages the overarching technical hiring pipeline that Gabe feeds into.\footnote{\href{https://theorg.com/org/google-x}{TheOrg: Google X Organizational Chart}}
    \item \textbf{Jackie Ruffner:} People Operations Program Manager at X.\footnote{\href{https://theorg.com/org/google-x}{TheOrg: Google X Organizational Chart}}
\end{itemize}

\section{Priority 2: The Core ML Pod (Your Future Team)}
When Gabe asks about your technical fit, you must align your background with these specific individuals. This is a heavy-science pod needing an engineering anchor.
\begin{itemize}[leftmargin=*]
    \item \textbf{Ali Ahmadalipour:} Applied Research Scientist (Geospatial \& Agentic AI, Climate ML).\footnote{\href{https://www.linkedin.com/in/ahmadalipour/}{Ali Ahmadalipour - LinkedIn Profile}}
    \item \textbf{Caroline Jaffe:} Geospatial AI and Climate Researcher (Ph.D., MIT).\footnote{\href{https://www.linkedin.com/in/carolinejaffe/}{Caroline Jaffe - LinkedIn Profile}}
    \item \textbf{Behzad Vahedi:} Spatial Machine Learning \& Data Science (Ph.D.).\footnote{\href{https://www.linkedin.com/in/behzad-vahedi/}{Behzad Vahedi - LinkedIn Profile}}
    \item \textbf{Marie Pelagie Elimbi Moudio, Ph.D. (Departed MLOps Engineer):} An early founding contributor to the MLOps pipeline who left in late 2023 for an Applied Scientist role at Amazon. Her departure, alongside Charles, confirms the pod currently has zero senior deployment/ops engineers.\footnote{\href{https://www.linkedin.com/in/marie-pelagie-elimbi-moudio-481335105/}{Marie Pelagie Elimbi Moudio - LinkedIn Profile}}
    \item \textbf{Nicholas Aflitto, Ph.D.:} Lead Ecologist at Google X and Assistant Professor of Forest Ecosystem Health at UVM. He bridges the gap between deep learning and ecological interactions (e.g., how ecosystems respond to wildfires).\footnote{\href{https://nickaflitto.com/}{Nicholas Aflitto - Ecological Research Profile}}
    \item \textbf{Sebastian Pilarski, Ph.D.:} AI Resident and Research Scientist at Google. Co-author on both AlphaEarth and Earth AI papers.\footnote{\href{https://sebastianpilarski.com/}{Sebastian Pilarski - AI Residency Profile}}
    \item \textbf{Luc Houriez \& Teo Honda Scully:} AI Researchers / Stanford PhD Candidates heavily involved in the AlphaEarth data generation pipelines.\footnote{Stanford Profiles / arXiv Co-Authorship.}
    \item \textbf{Charles Spirakis (Departed Former Head of Engineering):} Served as Head of Engineering for Bellwether from 2021 to 2024 before leaving to become CTO at Autoscience Institute. \textbf{This is highly critical:} You are not just joining the team, you are interviewing to directly backfill his leadership position and scale the V2 architecture.\footnote{\href{https://www.linkedin.com/in/charles-spirakis/}{Charles Spirakis - LinkedIn Profile}}
\end{itemize}

\section{Priority 3: Bellwether Leadership \& Mandate (The Business Case)}
If they ask "Why Bellwether?", you must understand their specific 2024--2026 trajectory and the recent leadership vacuum.
\begin{itemize}[leftmargin=*]
    \item \textbf{The Mandate:} A prediction engine for planetary-scale environmental changes, focusing on 5-year wildfire risk forecasting and disaster response. Named a \textbf{TIME Best Invention of 2024}.\footnote{\href{https://time.com/7094770/bellwether/}{TIME Best Inventions 2024: Bellwether}}
    \item \textbf{The ``Phase 2" Leadership Vacuum:} Both the original Founder/GM (Sarah Russell, departed Sept 2025) and the original Head of Engineering (Charles Spirakis, departed 2024) have exited the project. This strongly signals that Bellwether is transitioning from a "V1 prototype" built by the founders into a "V2 commercial scaling" phase requiring new, hardened technical leadership (you).
    \item \textbf{Rich Mazzola (Partnerships \& GTM):} The business anchor. He deploys the engine's forecasts to institutional partners. Under his mandate, Bellwether has established three massive verticals:\footnote{\href{https://www.linkedin.com/in/richamazzola/}{Rich Mazzola - LinkedIn Profile}}
    \begin{itemize}
        \item \textbf{DoD / Defense:} Officially partnered with the \textbf{Defense Innovation Unit (DIU)} and \textbf{National Guard} to automate aerial damage assessments using AlphaEarth.\footnote{\href{https://x.company/projects/bellwether/}{Project Bellwether - X, the moonshot factory (Official Site)}}
        \item \textbf{Federal / State Emergency:} Collaborating with the \textbf{U.S. Forest Service} on fire-spread simulation and deployed by state agencies for \textbf{Hurricane Helene} recovery.\footnote{Public deployments verified via news releases (e.g., Forbes, SmartEnergyDecisions).}
        \item \textbf{Enterprise / B2B SaaS:} Partnered with \textbf{SwissRe Reinsurance Solutions} to integrate AlphaEarth forecasting into the global insurance market.
    \end{itemize}
\end{itemize}

\section{Priority 4: X Executive Oversight (Late-Stage Interviews)}
These Directors evaluate Bellwether from a high level to decide if it graduates into an independent Alphabet company. You will likely interview with them in final rounds to prove your commercial maturity.
\begin{itemize}[leftmargin=*]
    \item \textbf{Astro Teller:} ``Captain of Moonshots" (CEO of X).\footnote{\href{https://www.astroteller.net/}{Astro Teller - Official Profile}}
    \item \textbf{John ``Ivo" Stivoric:} Vice President at X. Oversees the broader climate portfolio.\footnote{Public profile citations via X, the Moonshot Factory.}
    \item \textbf{David Andre:} Chief Science Officer at X. Provides scientific oversight across moonshots.\footnote{\href{https://www.linkedin.com/in/davidandre/}{David Andre - LinkedIn Profile}}
    \item \textbf{Rachel Payne / Mayank Girdhar / Jon Gold:} Managing Directors at X focused on commercialization, climate policy, and spin-outs.\footnote{\href{https://theorg.com/org/google-x}{TheOrg: Google X Organizational Chart} / Deployed Together Profile}
    \item \textbf{Josh Jeffery:} Director, Google X Strategy and Business Operations.\footnote{\href{https://www.linkedin.com/in/joshuajeffery/}{Josh Jeffery - LinkedIn Profile}}
\end{itemize}

\section{Priority 5: Parallel Sister Moonshots (Broader Context)}
X operates highly siloed projects, but engineers cross-pollinate. Knowing the sister grid-decarbonization moonshot (Tapestry) shows deep insider awareness.
\begin{itemize}[leftmargin=*]
    \item \textbf{Project Tapestry Leadership:} Page Crahan (GM), Prune Chabert (Product Manager), Charlotte Park (Head of People).\footnote{CERAWeek Profile / Volts Podcast / From Day One.}
    \item \textbf{Circularity Moonshot Leadership:} Rey Banatao (Director \& Project Lead).\footnote{\href{https://citris-uc.org/}{CITRIS Researcher Profile: Rey Banatao}}
\end{itemize}

\section{Priority 6: Technical Footprint \& Core Publications}
To win the technical screen, you must understand their proprietary stack and their recent academic releases.

\subsection*{Open Source, APIs, and Repositories}
\textbf{Status: Strictly Proprietary.} Project Bellwether does \textit{not} maintain public GitHub repositories, open-source libraries, or public-facing APIs. Their system is an enterprise-to-government prediction engine. In a recent Harvard Business School podcast, former GM Sarah Russell explicitly addressed the uncertainty of their code strategy, stating: ``Who will go open source? I don't know." \textbf{Takeaway:} Do not ask for API keys or GitHub links; instead, ask how they handle internal model versioning and massive proprietary data ingestion from partners like the National Guard.

\subsection*{Academic Publications (The Tech Stack)}
The core ML pod is actively publishing their architectural breakthroughs using DeepMind's infrastructure. \textbf{You should understand these 2025 arXiv preprints immediately before the call:}
\begin{itemize}[leftmargin=*]
    \item \textbf{``Scalable Geospatial Data Generation Using AlphaEarth Foundations Model" (arXiv:2507.20507, 2025):} Confirms Bellwether utilizes Google DeepMind's \textit{AlphaEarth} foundation model. \textbf{What it is:} AlphaEarth processes multi-modal data (satellite, radar, LiDAR) into 64-dimensional vectors embedding precise 10-meter Earth squares. It acts as a "virtual satellite," allowing Bellwether to interpolate missing environmental data via Space Time Precision.
    \item \textbf{``Earth AI: Unlocking Geospatial Insights with Foundation Models and Cross-Modal Reasoning" (arXiv:2510.18318, 2025):} Confirms Bellwether is tied into the broader "Google Earth AI" initiative. \textbf{What it is:} A \textit{Gemini-powered Geospatial Reasoning Agent} that takes natural language queries (e.g., "find all flooded roads"), orchestrates multiple underlying planet-scale models simultaneously, and fuses the multimodal results into actionable output.
\end{itemize}

\section{Priority 7: Your Core Pitch (Mapping your Resume to their Science)}
When Gabe asks why you are a fit for the Lead role, directly map your battle-tested infrastructure experience to the specific scientific breakthroughs the pod is currently publishing:

\begin{itemize}[leftmargin=*]
    \item \textbf{Aligning with Ali (Climate Extremes \& Foundation Models):} Tell Gabe you know Ali's work transitioning from flood damage prediction to Agentic AI (AlphaEarth/Gemini). Pitch that you developed real-time 1m flood inundation models at \textbf{TACC} for Texas Emergency Management, and you are actively building spatial Agentic RAG workflows at \textbf{PoliBOM}. You can bridge his science to production.
    \item \textbf{Pitching the MLOps/Wildfire Vacuum (Marie's old domain):} Marie built the early MLOps and wildfire hazard disruption models before leaving for Amazon. Pitch that you can step in to take ownership of the abandoned landscape optimization pipelines she left behind, leveraging your work spinning out \textbf{Summit Geospatial} to build the high-resolution, seamless state-wide LiDAR DEM mosaics required to feed those hazard models.
    \item \textbf{Aligning with Behzad (Spatial Deep Learning):} Behzad's core technical focus is spatial deep learning and custom classification layers (Focal Loss). Pitch your work at \textbf{Homecastr}, where you engineered a nationwide diffusion architecture using custom \textit{spatial inducing tokens} and per-horizon loss normalization, proving you know how to operate and train complex spatial neural networks at scale.
    \item \textbf{Aligning with the Business Mandate (DIU \& National Guard):} Bellwether deploys these models to the DoD and federal agencies. Hammer home that you competed in \textbf{DARPA World Modelers} (scaling 800,000 node distributed jobs) and partnered directly with the \textbf{US Army Corps of Engineers} and \textbf{NOAA} to deliver disaster-resilience tools. You already speak their customers' language perfectly.
\end{itemize}

\end{document}
